% LaTeX companion of 07-orthogonal-transformations.typ. Both formats are editable.
\chapter{Orthogonal transformations}

\section{WS 18, p. 1: definition and matrix criteria}

\begin{definition}{\kn{Orthogonal transformation}}

\(T:\mathbb{R}^{n}\rightarrow\mathbb{R}^{n}\) is orthogonal (正交变换) if it preserves dot products: \(\forall x,y \in \mathbb{R}^{n}\), \(x \cdot y = T(x) \cdot T(y)\).

\end{definition}

\begin{theorem}{Equivalent characterizations}

\(T:\mathbb{R}^{n}\rightarrow\mathbb{R}^{n}\) is orthogonal iff it preserves the length of vectors: \(\forall x \in \mathbb{R}^{n}\), \(\left\| {T(x)} \right\| = \left\| x \right\|\).

\end{theorem}

因而 linear trans 保留 dot product iff 保留 length. The proof expands \(\left\| {T\left( {x + y} \right)} \right\|^{2} = \left\| {x + y} \right\|^{2}\):

\[\left\| {T(x)} \right\|^{2} + 2T(x) \cdot T(y) + \left\| {T(y)} \right\|^{2} = \left\| x \right\|^{2} + 2x \cdot y + \left\| y \right\|^{2},\]

then cancels equal norm terms.

(a) An orthogonal trans \(T\) is injective: \(T(x) = 0\rightarrow\left\| {T(x)} \right\| = 0\rightarrow\left\| x \right\| = 0\), so \(\text{ker}\ T = \left\{ 0 \right\}\) and \(T\) is inj. (b) It is an isomorphism because a same-dimension linear trans is inj iff surj. (c) The standard matrix of \(T\) has orthonormal columns. (d) The composition of orthogonal trans is orthogonal, since

\[\left\| {T_{k} \cdot T_{k - 1} \cdot \ldots \cdot T_{1}(x)} \right\| = \ldots = \left\| x \right\|.\]

\begin{definition}{Orthogonal matrix}

A square matrix \(A\) is orthogonal if \(A^{t}A = I_{n}\) (即 \(A^{t} = A^{- 1}\)).

\end{definition}

If \(A = \begin{pmatrix}
v_{1} & \ldots & v_{n}
\end{pmatrix}\), then the \(ij\)th entry of \(A^{t}A\) is \(v_{i} \cdot v_{j}\). Therefore \(A\) is orthogonal iff its cols are orthonormal. 我们也由此知道：由 orthonormal basis 组成的 matrix \(A\)， \(A^{- 1}\) 就是 \(A^{t}\).

The worksheet proves \(\left( {AB} \right)^{t} = B^{t}A^{t}\): the \(ij\)th entry of \(AB\) is the \(i\)th row of \(A\) dot the \(j\)th col of \(B\), while the \(ij\)th entry of \(B^{t}A^{t}\) is \(b_{i} \cdot a_{j}\). Thus they agree.

\section{WS 18, p. 2: products and change of basis}

If \(A\) is orthogonal, then \(A^{- 1}\) (也是 \(A^{t}\)) is orthogonal, since \(\mathbb{A}^{t} = I_{n}\) implies \(\left( A^{t} \right)\left( A^{t} \right)^{t} = \mathbb{A}^{t} = I_{n}\). 结论：这意味着对于 \(A\) 的 cols 是 orthonormal 的，那 \(A\) 的 rows 也是 orthonormal 的.

\begin{theorem}{Products of orthogonal matrices}

If \(A,B\) are orthogonal \(n \times n\) matrices, then \(AB\) is orthogonal.

\end{theorem}

Indeed,

\[\left( {AB} \right)\left( {AB} \right)^{t} = A\left( {BB^{t}} \right)A^{t} = I_{n},\]

so \(\left( {AB} \right)^{t} = \left( {AB} \right)^{- 1}\).

\begin{theorem}{Orthogonal maps and their matrices}

\(T:\mathbb{R}^{n}\rightarrow\mathbb{R}^{n}\) is orthogonal iff \(\lbrack T\rbrack_{\varepsilon}\) is orthogonal, where \(\varepsilon\) is the standard basis. More generally, \(T\) is orthogonal iff \(\lbrack T\rbrack_{\beta}\) is orthogonal for any orthonormal basis \(\beta\).

\end{theorem}

For the reverse direction,

\[T(a) \cdot T(b) = \left( {\lbrack T\rbrack_{\varepsilon}a} \right)\dot{\lbrack T\rbrack_{\varepsilon}b} = a_{\varepsilon_{\varepsilon}^{t{\lbrack T\rbrack}}}^{t{\lbrack T\rbrack}}b = a \cdot b.\]

Claim 1: If \(A,B\) are orthonormal basis matrices, their change-of-basis matrices are orthogonal. On WS 16, the entries are

\[S_{A\rightarrow B} = \begin{pmatrix}
{a_{1} \cdot b_{1}} & {a_{2} \cdot b_{1}} & \ldots & {a_{n} \cdot b_{1}} \\
{a_{1} \cdot b_{2}} & \ldots & & \\
\ldots & \ldots & {a_{n} \cdot b_{n}} &
\end{pmatrix}\]

and

\[S_{B\rightarrow A} = \begin{pmatrix}
{b_{1} \cdot a_{1}} & {b_{2} \cdot a_{1}} & \ldots \\
\ldots & \ldots & {b_{n} \cdot a_{n}}
\end{pmatrix} = S_{A\rightarrow B}^{t}.\]

Thus \(S_{A\rightarrow B} = S_{B\rightarrow A}^{- 1}\).

\section{WS 18, p. 3: conclusion}

Claim 2: orthogonal matrices 的 product 也是 orthogonal matrix. 这是因为 orthogonal transformations 的 composition 也是 orthogonal transformation， 且它的自身的 matrix 代表为 orthogonal matrix。

Claim 3: 如果 \(T\) orthogonal，则 \(\lbrack T\rbrack_{\beta}\) orthogonal for 任意 orthonormal basis \(\beta\):

\[\lbrack T\rbrack_{\beta} = S_{\varepsilon\rightarrow\beta}\lbrack T\rbrack_{\varepsilon}S_{\beta\rightarrow\varepsilon}.\]

All three factors are orthogonal, hence so is \(\lbrack T\rbrack_{\beta}\).

Claim 4: \(\beta\) 为任意 orthonormal basis; 如果 \(\lbrack T\rbrack_{\beta}\) orthogonal，则 \(T\) orthogonal. Since

\[\lbrack T\rbrack_{\beta} = S_{\varepsilon\rightarrow\beta}\lbrack T\rbrack_{\varepsilon}S_{\beta\rightarrow\varepsilon},\]

we obtain

\[\lbrack T\rbrack_{\varepsilon} = S_{\beta\rightarrow\varepsilon}\lbrack T\rbrack_{\beta}S_{\varepsilon\rightarrow\beta},\]

a product of orthogonal matrices; hence \(\lbrack T\rbrack_{\varepsilon}\) is orthogonal and therefore \(T\) is orthogonal.

总结：\(T\) is orthogonal (保留 dot product) \(= T\) 保留 length \(= T\) 保留 distance \(= T\) 把 \(\mathbb{R}^{n}\) 的某个 orthonormal basis map 到另一个 orthonormal basis \(= \lbrack T\rbrack_{\beta}\) 为 orthogonal 的，\(\beta\) 为任意 orthonormal basis \(= \lbrack T\rbrack_{\beta}\) 的 rows/cols 为一个 orthonormal basis of \(\mathbb{R}^{n}\).
